\documentclass[10pt]{article}

\usepackage{amsmath,amssymb,amsthm}
\usepackage{url}

\newtheorem{definition}{Definition}
\newtheorem{problem}{Problem}
\newtheorem{theorem}{Theorem}
\newtheorem{remark}[theorem]{Remark}

\newcommand{\lcs}{\operatorname{lcs}}
\newcommand{\Zh}{\widehat{\mathbb{Z}}}
\newcommand{\stick}{\boxdot}
\newcommand{\red}{\operatorname{red}}
\newcommand{\repl}{\operatorname{repl}}
\newcommand{\dom}{\operatorname{dom}}
\newcommand{\ran}{\operatorname{ran}}
\newcommand{\BSP}{\text{BSP}}

\title{Doubly-Periodic LCS in the BSP Model - Easy}
\author{Nikita Gaevoy}
\date{}

\begin{document}

\maketitle

\section{Preliminaries}

The whole of this section is a summary of published work: Subsections
\ref{sec:strings}--\ref{sec:sequential} follow Gaevoy, Zolotov and
Tiskin~\cite{GZT25} and, through it, Tiskin's
framework~\cite{Tiskin08,TiskinArXiv,Tiskin15,Tiskin09}; Subsections
\ref{sec:bsp}--\ref{sec:bspsticky} follow Valiant~\cite{Valiant90} and
Krusche and Tiskin~\cite{KT07,KT10,KruscheThesis}. Nothing in it is new.

\subsection{Strings, repeats and the LCS score}\label{sec:strings}

We work with strings over an arbitrary alphabet. Given a string $a$, we write
$a\langle i:j\rangle$ for its substring of length $j-i$, and we distinguish
between contiguous \emph{substrings} and not necessarily contiguous
\emph{subsequences}.

\begin{definition}[LCS score, {\cite[Definition 1]{GZT25}}]
The \emph{longest common subsequence} (LCS) score for a pair of strings,
denoted $\lcs(a,b)$, is the length of the longest string that is a
subsequence of both $a$ and $b$. Given strings $a,b$, the \emph{LCS problem}
asks for their LCS score.
\end{definition}

\begin{definition}[Repeats, {\cite[Definition 2]{GZT25}}]
A string of the form $a^k=\text{``}aa\ldots a\text{''}$ ($k$ times) is called
a \emph{$k$-repeat}, or simply a \emph{repeat}. More generally, a string of
the form $a^ka'$, where $a'$ is a prefix of $a$, is called \emph{periodic}. A
doubly infinite string $a^\infty=\text{``}\ldots aaaa\ldots\text{''}$ is
called an \emph{infinite repeat}.
\end{definition}

Throughout, $a$ and $b$ are strings of length $m$ and $n$ respectively, and
$k,\ell\ge 1$ are integers. The \emph{doubly-periodic LCS problem} asks for
the score $\lcs(a^k,b^\ell)$, given $a$, $b$, $k$ and $\ell$; the input size
is thus $O(m+n+\log k+\log\ell)$, while the two compared strings have lengths
$mk$ and $n\ell$.

Half-integers are denoted by hatted letters ($\hat\imath$, $\hat\jmath$,
\ldots), and the set of all half-integers is denoted $\Zh$. We write
$\langle m:n\rangle$ for the set of half-integers strictly between the
integers $m,n$, and $[m:n]$ for the usual integer interval. Matrix elements
are written $M\langle\hat\imath;\hat\jmath\rangle$ when indexed by
half-integers and $M[i;j]$ when indexed by integers.

\subsection{Semi-local LCS and its kernel}

\begin{definition}[Semi-local LCS, {\cite[Definition 3]{GZT25}}]
Given strings $a,b$, the \emph{semi-local LCS} (SLCS) problem asks for the LCS
scores of the whole of $a$ against every substring of $b$ (string-substring
LCS), of every prefix of $a$ against every suffix of $b$ (prefix-suffix LCS),
and of the two symmetric variants.
\end{definition}

\begin{definition}[LCS grid, {\cite[Definition 4]{GZT25}}]\label{def:grid}
The \emph{LCS grid} $G_{a,b}[0:m;0:n]$ for strings $a,b$ is a directed acyclic
graph on the node set $[0:m;0:n]$ containing, for all $i\in[0:m]$,
$\hat\imath\in\langle 0:m\rangle$, $j\in[0:n]$,
$\hat\jmath\in\langle 0:n\rangle$, the horizontal edge
$[i;\hat\jmath^-]\to[i;\hat\jmath^+]$ and the vertical edge
$[\hat\imath^-;j]\to[\hat\imath^+;j]$; and, whenever
$a\langle\hat\imath\rangle$ matches $b\langle\hat\jmath\rangle$, the diagonal
\emph{match edge} $[\hat\imath^-;\hat\jmath^-]\to[\hat\imath^+;\hat\jmath^+]$.
A cell carrying a match edge is a \emph{match cell}, and otherwise a
\emph{mismatch cell}.
\end{definition}

Padding $b$ with $m$ wildcard characters on either side, so that
$b_{\mathrm{pad}}\langle-m:m+n\rangle=\text{?}^m b\,\text{?}^m$, the
semi-local scores are collected in a single matrix.

\begin{definition}[SLCS matrix, {\cite[Definition 5]{GZT25}}]
The \emph{SLCS matrix} for strings $a,b$ is the matrix
$H_{a,b}[-m:n;0:m+n]$ with $H_{a,b}[i;j]=
\lcs\bigl(a,b_{\mathrm{pad}}\langle i:j\rangle\bigr)$ for $i\le j$ and
$H_{a,b}[i;j]=j-i$ for $i\ge j$.
\end{definition}

\begin{definition}[Dominance sums, {\cite[Definition 10]{GZT25}}]
Let $D\langle I;J\rangle$ be a matrix. Its \emph{dominance-sum matrix}
$D^\Sigma[I;J]$ is defined by
$D^\Sigma[i;j]=\sum D\langle i:;\,:j\rangle$ for all $i\in[I]$, $j\in[J]$.
\end{definition}

\begin{theorem}[SLCS matrix canonical representation,
{\cite[Theorem 33]{GZT25}}]\label{thm:canon}
The SLCS matrix $H_{a,b}$ is unit-anti-Monge:
\[
H_{a,b}[i;j]=-i+j+m-P^\Sigma_{a,b}[i;j]
\]
for all $i,j$, where $P_{a,b}$ is a permutation matrix.
\end{theorem}

\begin{definition}[SLCS kernel, {\cite[Definition 34]{GZT25}}]
The \emph{SLCS kernel} for strings $a,b$ is the permutation matrix
$P_{a,b}\langle-m:n;0:m+n\rangle$ determined by Theorem~\ref{thm:canon}. The
\emph{string-substring LCS subkernel} is the subpermutation matrix
$P^{\nearrow}_{a,b}=P_{a,b}\langle 0:;\,:n\rangle$.
\end{definition}

\subsection{Sticky multiplication}

\begin{definition}[Tropical and sticky multiplication,
{\cite[Definitions 11, 12]{GZT25}}]
For matrices $A[I;J]$, $B[J;K]$, the \emph{tropical} (that is, $(\min,+)$)
product is
$(A\odot B)[i;k]=\min_{j\in[J]}\bigl(A[i;j]+B[j;k]\bigr)$. For matrices
$D,E,F$, the \emph{implicit tropical}, or \emph{sticky}, product is defined by
\[
D\stick E=F
\qquad\text{iff}\qquad
D^\Sigma\odot E^\Sigma=F^\Sigma .
\]
\end{definition}

\begin{definition}[{\cite[Definition 13]{GZT25}}]
The \emph{sticky multiplication monoid of permutation matrices} $H_n$ is the
set of all $n\times n$ permutation matrices under sticky multiplication.
\end{definition}

The monoid $H_n$ is isomorphic to the \emph{sticky braid monoid} of degree
$n$, presented by generators $\tau_1,\ldots,\tau_{n-1}$ subject to
$\tau_i^2=\tau_i$, far commutativity and the braid
relations~\cite[Definition 15, Theorem 19]{GZT25}. The algorithmic content of
the framework is the following theorem.

\begin{theorem}[Tiskin--Sakai, {\cite{Tiskin15,Sakai11}},
{\cite[Theorem 14]{GZT25}}]\label{thm:tiskinsakai}
Let $P\stick Q=R$, where $P,Q,R$ are (sub)permutation matrices with at most
$n$ nonzeros each. Given the nonzeros of $P,Q$, the nonzeros of $R$ can be
obtained in time $O(n\log n)$.
\end{theorem}

\begin{theorem}[SLCS kernel composition, {\cite[Theorem 35]{GZT25}}]
Let $a=a'a''$, where $a',a''$ are of length $m',m''$. Then
$P^{\nearrow}_{a,b}=P^{\nearrow}_{a',b}\stick P^{\nearrow}_{a'',b}$, and the
product can be obtained in time $O(n\log N)$, where
$N=\min\{m',m'',n\}$.
\end{theorem}

\subsection{The affine extension}\label{sec:affine}

\begin{definition}[{\cite[Definition 20]{GZT25}}]\label{def:affineperm}
A \emph{skew-affine permutation} of degree $n$ and \emph{skew} $s$ is a
bijection $\tilde P:\Zh\to\Zh$ such that $\tilde P(\hat\imath+n)=
\tilde P(\hat\imath)+n$ for all $\hat\imath\in\Zh$, and
$\sum_{\hat\imath\in\langle 0:n\rangle}\tilde P(\hat\imath)
=\sum_{\hat\imath\in\langle 0:n\rangle}\hat\imath+s$. A skew-affine
permutation with skew $0$ is called \emph{affine}.
\end{definition}

Permutation matrices, dominance sums, tropical multiplication and sticky
multiplication all extend from the finite to the affine
case~\cite{Pflueger22}, yielding the monoid $\tilde H_n$ of affine
permutation matrices of degree $n$ under sticky multiplication, isomorphic to
the affine sticky braid monoid of degree
$n$~\cite[Definitions 22, 24, Theorem 25]{GZT25}. Sticky multiplication is
well defined on skew-affine permutations as well: the skew of a product is
the sum of the skews of the multiplicands, and applying the skews commutes
with taking the product~\cite[Remark 23]{GZT25}. Throughout, an affine
permutation is represented by the images of a single period, that is, by the
restriction of $\tilde P$ to $\langle 0:n\rangle$, taken as a mapping into
$\Zh$~\cite[Section 4.1]{GZT25}.

\begin{definition}[{\cite[Definitions 26, 28]{GZT25}}]\label{def:phigamma}
Let $f:\langle 0:n\rangle\to\Zh$ be an injection whose $n$ images are pairwise
non-congruent modulo $n$. The \emph{affine replication} $\repl_n f$ is the
unique affine permutation of degree $n$ whose restriction to
$\langle 0:n\rangle$ equals $f$. An affine permutation $\tilde\Phi$ is
\emph{finite-type} if $\tilde\Phi|_{\langle 0:N\rangle}$ is a finite
permutation; an affine permutation $\tilde\Gamma$ is \emph{Grassmannian} if
$\tilde\Gamma|_{\langle 0:N\rangle}$ is monotone increasing.
\end{definition}

\begin{theorem}[{\cite[Proposition 29]{GZT25}}]\label{thm:decomp}
Any affine permutation $\tilde P$ can be represented as a product
$\tilde P=\tilde\Phi\cdot\tilde\Gamma$ with $\tilde\Phi$ finite-type and
$\tilde\Gamma$ Grassmannian. Such a \emph{$\Phi\Gamma$-decomposition} can be
computed in time $O(N\log N)$ and linear space; the symmetric
\emph{$\Gamma\Phi$-decomposition}
$\tilde P=\tilde\Gamma'^{\,-1}\cdot\tilde\Phi'$ is computed in the same time.
\end{theorem}

The proof of Theorem~\ref{thm:decomp} consists of sorting the elements of
$\langle 0:N\rangle$ by the order of their $\tilde P$-images, and one
composition of affine permutations; the $O(N\log N)$ term is the sort.

Given strings $a$ of length $m$ and $b$ of length $n$, the string-substring
component of the semi-local problem for the pair $a$, $b^\infty$ again admits
a kernel, denoted $P_{a,b^\infty}$. Provided that every character of $a$
occurs in $b$ at least once --- which can be assumed without loss of
generality, since a character of $a$ absent from $b$ is never matched --- the
kernel $P_{a,b^\infty}$ is a skew-affine permutation of degree $n$ and skew
$m$~\cite[Section 3.3]{GZT25}.

\begin{theorem}[Affine kernel composition,
{\cite[Theorem 40]{GZT25}}]\label{thm:affinecompose}
Let $a=a'a''$. Then
$P_{a,b^\infty}=P_{a',b^\infty}\stick P_{a'',b^\infty}$, the sticky product
being taken on skew-affine permutations as above.
\end{theorem}

\begin{theorem}[Wraparound combing, \cite{Tiskin09},
{\cite[Section 3.3]{GZT25}}]\label{thm:combing}
Given two strings $a,b$ of lengths $m,n$, the kernel $P_{a,b^\infty}$ can be
computed in time $O(mn)$ and space $O(n)$.
\end{theorem}

Wraparound combing initialises a saturated affine embedded braid on the
infinite periodic LCS grid and iterates over its cells: the outer loop runs
over the characters of $a$, north to south; the inner loop runs over one whole
period of $b^\infty$, beginning at a match cell in the current row, with an
entanglement check that respects the periodic structure of the
grid~\cite[Section 3.3]{GZT25}. The choice of the starting period may differ
between iterations of the outer loop, and the inner loop is a single
left-to-right sweep across a period.

\subsection{The sequential doubly-periodic algorithm}\label{sec:sequential}

\begin{theorem}[Affine sticky multiplication,
{\cite[Theorems 42, 44, Remark 43]{GZT25}}]\label{thm:affinesticky}
Let $\tilde P,\tilde Q$ be skew-affine permutations of degree $n$. The
following procedure computes $\tilde P\stick\tilde Q$, in time
$O(n\log n)$ and space $O(n)$.
\begin{enumerate}
\item Compute the $\Gamma\Phi_{3n}$-decomposition
$\tilde P=\tilde\Gamma\cdot\tilde\Phi$ and the
$\Phi\Gamma_{3n}$-decomposition $\tilde Q=\tilde\Phi'\cdot\tilde\Gamma'$. Put
$P=\tilde\Phi|_{\langle 0:3n\rangle}$,
$Q=\tilde\Phi'|_{\langle 0:3n\rangle}$.
\item Compute the finite sticky product $P\stick Q$, and from it the
relative reduction $\red(P,Q):=P^{-1}\cdot(P\stick Q)$.
\item Compute $S=\red(P,Q)|_{\langle n:2n\rangle}\cdot\tilde\Gamma'$, which is
an injection $\langle n:2n\rangle\to\Zh$; then
$\red(\tilde P,\tilde Q)=\repl_n S$.
\item Output $\tilde P\cdot\red(\tilde P,\tilde Q)$.
\end{enumerate}
\end{theorem}

Thus a single affine sticky multiplication of degree $n$ costs two
decompositions as in Theorem~\ref{thm:decomp} (that is, two sorts of $3n$
keys), one finite sticky multiplication of order $3n$ as in
Theorem~\ref{thm:tiskinsakai}, and a constant number of compositions and
inversions of permutations of order $O(n)$.

\begin{theorem}[{\cite[Theorem 45]{GZT25}}]\label{thm:kernelpower}
Given strings $a,b$ of lengths $m,n$, the kernel $P_{a^k,b^\infty}$ can be
computed in time $O(mn+n\log n\log k)$ and space $O(n)$.
\end{theorem}

The proof of Theorem~\ref{thm:kernelpower} is: build $P_{a,b^\infty}$ by
wraparound combing (Theorem~\ref{thm:combing}), then raise it to the power $k$
in $\tilde H_n$ by binary exponentiation, using
Theorem~\ref{thm:affinecompose} for correctness and
Theorem~\ref{thm:affinesticky} as the subroutine.

\begin{theorem}[{\cite[Theorem 46]{GZT25}}]\label{thm:score}
Given the kernel $P_{u,b^\infty}$ for a string $u$ and a string $b$ of length
$n$, the score $\lcs(u,b^\ell)$ can be computed in time $O(n)$ by
\[
\lcs\bigl(u,b^\ell\bigr)
=\sum_{0<\hat\jmath<n}
\min\left\{\ell,\ \left\lfloor
\frac{P_{u,b^\infty}(\hat\jmath)}{n}\right\rfloor\right\}.
\]
\end{theorem}

\begin{theorem}[{\cite[Corollary 47]{GZT25}}]\label{thm:sequential}
The score $\lcs(a^k,b^\ell)$ can be computed in time
$O(mn+n\log n\log k)$ and space $O(n+m)$.
\end{theorem}

\subsection{The BSP model}\label{sec:bsp}

\begin{definition}[BSP computer, \cite{Valiant90}]
A \emph{BSP computer} is described by three parameters $(p,g,l)$: the number
$p$ of processors, each with its own local memory; the \emph{communication
gap} $g$, describing the rate at which data can be transmitted continuously by
the network relative to the computation speed of a single processor; and the
\emph{latency} $l$, the overhead of starting up communication. A BSP
computation is divided into \emph{supersteps}, each consisting of a phase of
local computation followed by a communication phase, and terminated by a
barrier synchronisation of all processors.
\end{definition}

Following~\cite[Section 2.2]{KruscheThesis}, a BSP algorithm on input of size
$N$ is measured by four quantities: the \emph{local computation}
$W(N,p)$, the sum over all supersteps of the maximum number of local
operations performed by a processor in that superstep; the
\emph{communication} $H(N,p)$, the sum over all supersteps of the maximum
number of data units sent or received by a processor in that superstep; the
\emph{memory} $M(N,p)$, the maximum amount of local storage used by a
processor at any point; and the \emph{synchronisation} $S$, the number of
supersteps. The cost of the computation is $W+g\cdot H+l\cdot S$.

\begin{definition}[{\cite[Definitions 2.3.1, 2.3.3, 2.3.4, 2.3.5]{KruscheThesis}}]
\label{def:scalable}
Let the sequential reference algorithm run in time $\mathcal W(N)$ and space
$\mathcal M(N)$, and let $\mathcal I(N)$ be the maximum of the input and
output size. A BSP algorithm is \emph{work-optimal} if
$W(N,p)=O\bigl(\mathcal W(N)/p\bigr)$; it has \emph{scalable communication}
if $H(N,p)=O\bigl(\mathcal I(N)/p^c\bigr)$ for some constant $c>0$; it has
\emph{scalable memory} if $M(N,p)=O\bigl(\mathcal M(N)/p^c\bigr)$ for some
constant $c>0$; and it is \emph{synchronisation-efficient} if $S$ is not a
function of $N$.
\end{definition}

\begin{remark}\label{rem:slackness}
As usual in the BSP model, the bounds of Definition~\ref{def:scalable} are
attainable only under a \emph{slackness} condition relating $N$ and $p$; for
instance, BSP parallel prefix runs in $W=O(N/p)$, $H=O(p)$, $S=O(1)$ provided
$N\ge p^2$~\cite[Example 2.3.6]{KruscheThesis}. All the bounds stated below
carry such conditions explicitly.
\end{remark}

We use two standard BSP primitives. Sorting $N$ keys can be performed in
$W=O\bigl((N\log N)/p\bigr)$, $H=M=O(N/p)$ and $S=O(1)$ supersteps under
polynomial slackness~\cite{Goodrich99,GV94}. Applying a permutation of
$[0:N]$ that is distributed evenly across the processors --- equivalently,
composing or inverting two such permutations --- is a balanced $h$-relation
with $h=O(N/p)$, and thus costs $W=H=O(N/p)$ in $S=O(1)$ supersteps.

\subsection{Sticky multiplication in parallel}\label{sec:bspsticky}

Sticky multiplication of \emph{finite} permutation matrices has been
parallelised in the BSP model, and in more restrictive models. The following
two results are the ones relevant here.

\begin{theorem}[{\cite{KT07}}, {\cite[Section 4.2]{KruscheThesis}}]
\label{thm:kt07}
Given the nonzeros of two $n\times n$ permutation matrices $P_A$, $P_B$,
distributed equally across $p$ processors, the nonzeros of $P_C$ with
$P_C=P_A\stick P_B$ can be computed using
$W(n,p)=O\bigl(n^{1.5}/p\bigr)$,
$H(n,p)=M(n,p)=O\bigl(n/\sqrt p\bigr)$ and $S=O(1)$, under the slackness
condition $n>p^2$.
\end{theorem}

\begin{theorem}[{\cite{KT10}}, {\cite[Lemma 4.3.3]{KruscheThesis}}]
\label{thm:kt10}
Given the nonzeros of two $n\times n$ permutation matrices $P_A$, $P_B$,
distributed equally across $p<\sqrt[3]{n}$ processors, the nonzeros of $P_C$
with $P_C=P_A\stick P_B$ can be computed using
$W(n,p)=O\bigl((n\log n)/p\bigr)$, $H(n,p)=O\bigl((n/p)\log p\bigr)$,
$M(n,p)=O(n/p)$ and $S=O(\log p)$.
\end{theorem}

Theorem~\ref{thm:kt10} is work-optimal with respect to
Theorem~\ref{thm:tiskinsakai}, and has scalable communication within a
$\log p$ factor of $n/p$; Theorem~\ref{thm:kt07} is synchronisation-efficient
but is not work-optimal, and its communication and memory are scalable only
with $c=1/2$. Building on Theorem~\ref{thm:kt10}, Krusche and Tiskin obtain a
BSP algorithm for semi-local LCS of two strings of length $n$ running in
$W=O(n^2/p)$, $H=M=O(n/\sqrt p)$ and
$S=O(\log^2p)$~\cite[Theorem 4.4.1]{KruscheThesis}. Tiskin~\cite{Tiskin20}
subsequently improved the known upper bounds on communication and
synchronisation for parallel semi-local LCS, and gave the first constant-sync
work-optimal LCS algorithm. In the more restrictive massively parallel
computation (MPC) model of~\cite{KSV10}, Koo~\cite{Koo24} gave an
$O(1)$-round fully scalable deterministic algorithm for the same
multiplication.

\section{Problem formulation}

\begin{problem}[Doubly-periodic LCS in BSP]\label{prob:lcs}
Prove that there is a BSP algorithm which, given strings $a,b$ of lengths
$m,n$ distributed equally across $p$ processors, together with $k$ and
$\ell$, computes $\lcs(a^k,b^\ell)$ using
\begin{gather*}
W=O\!\left(\frac{mn+n\log n\log k}{p}\right),\qquad
H=M=O\!\left(n+\frac{m}{p}\right),
\\
S=O\bigl(\log^2p+\log p\log k\bigr),
\end{gather*}
under the slackness conditions $n>p^3$ and $m\ge p\log n$; that is, which is
work-optimal with respect to Theorem~\ref{thm:sequential}.
Better still, and the form the problem is really after, is an algorithm
meeting the same bounds on $W$ and $S$ with
\[
H=M=O\!\left(\frac{m}{p}+\frac{n}{p^c}\right)
\qquad\text{for some constant }c>0,
\]
possibly under a stronger slackness condition; that is, one with scalable
communication and scalable memory in the sense of
Definition~\ref{def:scalable}.
\end{problem}

\begin{remark}\label{rem:monotone}
The scalable form of Problem~\ref{prob:lcs} implies the plain one, since
$m/p+n/p^c=O(n+m/p)$; the plain one is a solution as it stands, and the
scalable one is the same problem carried further.
The intended solution calls a BSP algorithm for affine sticky multiplication
in $\tilde H_n$ $O(\log p+\log k)$ times. That multiplication is the only part
of the pipeline for which no algorithm in any parallel model is currently in
the literature, and it is of independent interest in its own right.
\end{remark}

\begin{remark}\label{rem:sync}
A variant that is equally interesting asks to trade work for synchronisation:
using a synchronisation-efficient multiplication such as
Theorem~\ref{thm:kt07} in place of Theorem~\ref{thm:kt10} should give
$S=O(\log p+\log k)$ at the cost of a larger $W$ and of communication and
memory scalable only with $c=1/2$. Since the binary exponentiation of
Theorem~\ref{thm:kernelpower} is inherently sequential in its $\log k$
rounds, $S=\Omega(\log k)$ for any algorithm of this shape, so
$S=O(\log p+\log k)$ is the natural target. The same question in the MPC
model, where Koo's $O(1)$-round fully scalable
multiplication~\cite{Koo24} is available in place of
Theorem~\ref{thm:kt07}, asks for $O(\log k)$ rounds.
\end{remark}

\begin{remark}\label{rem:asymmetry}
The bounds are asymmetric in the two inputs: $k$ enters through
$\log k$ and $\ell$ only through Theorem~\ref{thm:score}. This is inherited
from the sequential algorithm, where the roles of $a^k$ and $b^\ell$ are
genuinely different --- the first string is combed, the second one is folded
into the affine structure. Swapping the two inputs gives the bound with $m$
and $n$, and $k$ and $\ell$, interchanged, and one may of course run whichever
of the two is cheaper.
\end{remark}

\begin{remark}\label{rem:general}
Following the published version of~\cite{GZT25}, the problem is posed for a
pair of repeats $a^k$, $b^\ell$. Theorem~\ref{thm:score} is stated there for
a repeat $b^\ell$, with the case of a general periodic string $b^\ell b'$
deferred to a full version that is not available; a solution of the
problem above is not expected to depend on
which of the two variants is used, since the difference lies entirely in the
$O(n)$ score-retrieval phase.
\end{remark}

\begin{thebibliography}{GZT25}

\bibitem[Goo99]{Goodrich99}
Michael~T. Goodrich.
\newblock Communication-efficient parallel sorting.
\newblock {\em SIAM Journal on Computing}, 29(2):416--432, 1999.

\bibitem[GV94]{GV94}
Alexandros~V. Gerbessiotis and Leslie~G. Valiant.
\newblock Direct bulk-synchronous parallel algorithms.
\newblock {\em Journal of Parallel and Distributed Computing}, 22(2):251--267,
  1994.

\bibitem[GZT25]{GZT25}
Nikita Gaevoy, Boris Zolotov, and Alexander Tiskin.
\newblock Doubly-periodic string comparison.
\newblock In Paola Bonizzoni and Veli M{\"a}kinen, editors, {\em 36th Annual
  Symposium on Combinatorial Pattern Matching (CPM 2025)}, volume 331 of {\em
  LIPIcs}, pages 13:1--13:19. Schloss Dagstuhl -- Leibniz-Zentrum f\"{u}r
  Informatik, 2025.

\bibitem[Koo24]{Koo24}
Jaehyun Koo.
\newblock An optimal {MPC} algorithm for subunit-{M}onge matrix multiplication,
  with applications to {LIS}.
\newblock In {\em 36th ACM Symposium on Parallelism in Algorithms and
  Architectures (SPAA 2024)}, pages 145--154, 2024.

\bibitem[Kru10]{KruscheThesis}
Peter Krusche.
\newblock {\em Parallel String Alignments: Algorithms and Applications}.
\newblock PhD thesis, University of Warwick, 2010.
\newblock \url{https://pkrusche.netlify.app/files/semithesis.pdf}.

\bibitem[KSV10]{KSV10}
Howard Karloff, Siddharth Suri, and Sergei Vassilvitskii.
\newblock A model of computation for {M}ap{R}educe.
\newblock In {\em 2010 ACM-SIAM Symposium on Discrete Algorithms (SODA 2010)},
  pages 938--948, 2010.

\bibitem[KT07]{KT07}
Peter Krusche and Alexander Tiskin.
\newblock Efficient parallel string comparison.
\newblock In {\em Parallel Computing: Architectures, Algorithms and
  Applications (ParCo 2007)}, volume~38 of {\em NIC Series}, pages 193--200.
  John von Neumann Institute for Computing, 2007.

\bibitem[KT10]{KT10}
Peter Krusche and Alexander Tiskin.
\newblock New algorithms for efficient parallel string comparison.
\newblock In {\em 22nd Annual ACM Symposium on Parallelism in Algorithms and
  Architectures (SPAA 2010)}, pages 209--216, 2010.

\bibitem[Pfl22]{Pflueger22}
Nathan Pflueger.
\newblock An extended {D}emazure product on integer permutations via min-plus
  matrix multiplication, 2022.
\newblock arXiv:2206.14227.

\bibitem[Sak11]{Sakai11}
Yoshifumi Sakai.
\newblock A fast algorithm for multiplying min-sum permutations.
\newblock {\em Discrete Applied Mathematics}, 159(17):2175--2183, 2011.

\bibitem[Tis08]{Tiskin08}
Alexander Tiskin.
\newblock Semi-local string comparison: Algorithmic techniques and
  applications.
\newblock {\em Mathematics in Computer Science}, 1(4):571--603, 2008.

\bibitem[Tis09]{Tiskin09}
Alexander Tiskin.
\newblock Periodic string comparison.
\newblock In {\em 20th Annual Symposium on Combinatorial Pattern Matching (CPM
  2009)}, volume 5577 of {\em Lecture Notes in Computer Science}, pages
  193--206. Springer, 2009.

\bibitem[Tis13]{TiskinArXiv}
Alexander Tiskin.
\newblock Semi-local string comparison: algorithmic techniques and
  applications, 2013.
\newblock arXiv:0707.3619v21.

\bibitem[Tis15]{Tiskin15}
Alexander Tiskin.
\newblock Fast distance multiplication of unit-{M}onge matrices.
\newblock {\em Algorithmica}, 71(4):859--888, 2015.

\bibitem[Tis20]{Tiskin20}
Alexander Tiskin.
\newblock Communication vs synchronisation in parallel string comparison.
\newblock In {\em 32nd ACM Symposium on Parallelism in Algorithms and
  Architectures (SPAA 2020)}, pages 479--489, 2020.

\bibitem[Val90]{Valiant90}
Leslie~G. Valiant.
\newblock A bridging model for parallel computation.
\newblock {\em Communications of the ACM}, 33(8):103--111, 1990.

\end{thebibliography}

\end{document}
